\documentclass[12pt]{article}

\usepackage[utf8]{inputenc}
\usepackage{amsmath, amssymb, amsfonts, cancel}
\usepackage{graphicx}
\usepackage[letterpaper, left=1in, right=1in, top=1in, bottom=1in]{geometry}
\usepackage{booktabs}
\usepackage{multirow}
\usepackage{float}
\usepackage{caption}
\usepackage{colortbl}
\usepackage{xcolor}
\definecolor{paleYellow}{RGB}{255,255,180}
\usepackage{physics}
\usepackage{hyperref}
\usepackage[spanish]{babel}

\usepackage{tikz}

\usetikzlibrary{
    positioning,
    shapes.geometric,
    arrows.meta
}

\title{
Tarea 1: Subredes IPv4\\[0.5em]
\large Interacción de redes
}

\author{
Jorge Alberto Suarez Saldana\\
\small Matrícula: 355992\\
\small Profesor: Raymundo Antonio González Grimaldo
}

\date{\today}

\begin{document}
\maketitle

\section{Topología}

\begin{figure}[ht]
    \centering

    \begin{tikzpicture}[
        router/.style={
            draw,
            thick,
            rounded corners,
            minimum width=1.5cm,
            minimum height=0.8cm,
            fill=gray!15
        },
        switch/.style={
            draw,
            thick,
            rectangle,
            minimum width=1.4cm,
            minimum height=0.7cm,
            fill=blue!10
        },
        server/.style={
            draw,
            thick,
            rectangle,
            minimum width=1.5cm,
            minimum height=0.7cm,
            fill=green!10
        },
        pc/.style={
            draw,
            thick,
            rectangle,
            minimum width=1.2cm,
            minimum height=0.7cm,
            fill=orange!10
        },
        subnet/.style={
            midway,
            fill=white,
            inner sep=2pt,
            font=\small
        },
        link/.style={
            thick
        }
    ]

    % =========================================================
    % ROUTERS
    % =========================================================

    \node[router] (R1) at (0,4) {R1};
    \node[router] (R2) at (6,4) {R2};
    \node[router] (R3) at (3,0) {R3};

    % =========================================================
    % ENLACES ENTRE ROUTERS
    % =========================================================

    % Subred 4
    \draw[link]
        (R1) -- node[subnet, above] {Subred 4}
        (R2);

    % Subred 5
    \draw[link]
        (R1) -- node[subnet, left] {Subred 5}
        (R3);

    % Subred 6
    \draw[link]
        (R2) -- node[subnet, right] {Subred 6}
        (R3);

    % =========================================================
    % SUBRED 1
    % =========================================================

    \node[switch] (S1) at (-3,4) {S1};

    \draw[link]
        (R1) -- node[subnet, above] {Subred 1}
        (S1);

    \node[server] (Serv1) at (-4.5,2.3) {Serv1};
    \node[server] (Serv2) at (-1.5,2.3) {Serv2};

    \draw[link] (S1) -- (Serv1);
    \draw[link] (S1) -- (Serv2);

    % =========================================================
    % SUBRED 2
    % =========================================================

    \node[switch] (S2) at (3,-2) {S2};

    \draw[link]
        (R3) -- node[subnet, right] {Subred 2}
        (S2);

    \node[pc] (PCA) at (1,-4) {PCA};

    \node[switch] (S3) at (5,-4) {S3};

    \draw[link] (S2) -- (PCA);
    \draw[link] (S2) -- (S3);

    \node[pc] (PCB) at (5,-6) {PCB};

    \draw[link] (S3) -- (PCB);

    % =========================================================
    % SUBRED 3
    % =========================================================

    \node[switch] (S4) at (9,4) {S4};

    \draw[link]
        (R2) -- node[subnet, above] {Subred 3}
        (S4);

    \node[server] (Serv3) at (9,2) {Serv3};

    \draw[link] (S4) -- (Serv3);

    \end{tikzpicture}

    \caption{Topología de la red y distribución de las subredes.}
    \label{fig:topologia-red}

\end{figure}

La red principal es:

\[
140.200.128.0/18
\]

Se requieren al menos \(8=2^3\) subredes, por lo que se toman 3 bits de la parte de hosts:

\[
/18+3=/21
\]

La nueva máscara es:

\[
255.255.248.0
\]

Cada subred contiene:

\[
2^{32-21}=2^{11}=2048
\]

direcciones totales, de las cuales 2046 son utilizables para hosts:

\[
2^{11}-2=2046
\]

\subsection{Subredes obtenidas}

\begin{table}[h]
\centering
\begin{tabular}{c|c|c|c|c}
\textbf{Subred} & \textbf{ID de red} & \textbf{1er host} &
\textbf{Último host} & \textbf{Broadcast}\\
\hline
1 & 140.200.128.0 & 140.200.128.1 & 140.200.135.254 & 140.200.135.255\\
2 & 140.200.136.0 & 140.200.136.1 & 140.200.143.254 & 140.200.143.255\\
3 & 140.200.144.0 & 140.200.144.1 & 140.200.151.254 & 140.200.151.255\\
4 & 140.200.152.0 & 140.200.152.1 & 140.200.159.254 & 140.200.159.255\\
5 & 140.200.160.0 & 140.200.160.1 & 140.200.167.254 & 140.200.167.255\\
6 & 140.200.168.0 & 140.200.168.1 & 140.200.175.254 & 140.200.175.255\\
7 & 140.200.176.0 & 140.200.176.1 & 140.200.183.254 & 140.200.183.255\\
8 & 140.200.184.0 & 140.200.184.1 & 140.200.191.254 & 140.200.191.255
\end{tabular}
\caption{División de la red principal en 8 subredes /21.}
\end{table}

\subsection{Asignación de subredes a la topología}

\begin{table}[h]
\centering
\begin{tabular}{c|l}
\textbf{Subred} & \textbf{Segmento}\\
\hline
1 & R1 -- S1 -- Serv1, Serv2\\
2 & R3 -- S2 -- PCA -- S3 -- PCB\\
3 & R2 -- S4 -- Serv3\\
4 & R1 -- R2\\
5 & R3 -- R1\\
6 & R2 -- R3\\
7 & Disponible\\
8 & Disponible
\end{tabular}
\caption{Asignación de las subredes a la topología.}
\end{table}

\subsection{Tabla de direccionamiento}

Se utiliza la primera dirección válida de cada LAN como puerta de enlace.
En los enlaces entre routers se utilizan las dos primeras direcciones válidas.

\begin{table}[h]
\centering
\small
\begin{tabular}{l|l|l|c|l}
\textbf{Dispositivo} & \textbf{Interfaz} & \textbf{IP} &
\textbf{Máscara/Prefijo} & \textbf{Gateway}\\
\hline
R1 & LAN (Subred 1) & 140.200.128.1 & 255.255.248.0 /21 & --\\
Serv1 & NIC & 140.200.128.2 & 255.255.248.0 /21 & 140.200.128.1\\
Serv2 & NIC & 140.200.128.3 & 255.255.248.0 /21 & 140.200.128.1\\
\hline
R3 & LAN (Subred 2) & 140.200.136.1 & 255.255.248.0 /21 & --\\
PCA & NIC & 140.200.136.2 & 255.255.248.0 /21 & 140.200.136.1\\
PCB & NIC & 140.200.136.3 & 255.255.248.0 /21 & 140.200.136.1\\
\hline
R2 & LAN (Subred 3) & 140.200.144.1 & 255.255.248.0 /21 & --\\
Serv3 & NIC & 140.200.144.2 & 255.255.248.0 /21 & 140.200.144.1\\
\hline
R1 & R1--R2 (Subred 4) & 140.200.152.1 & 255.255.248.0 /21 & --\\
R2 & R2--R1 (Subred 4) & 140.200.152.2 & 255.255.248.0 /21 & --\\
\hline
R3 & R3--R1 (Subred 5) & 140.200.160.1 & 255.255.248.0 /21 & --\\
R1 & R1--R3 (Subred 5) & 140.200.160.2 & 255.255.248.0 /21 & --\\
\hline
R2 & R2--R3 (Subred 6) & 140.200.168.1 & 255.255.248.0 /21 & --\\
R3 & R3--R2 (Subred 6) & 140.200.168.2 & 255.255.248.0 /21 & --
\end{tabular}
\caption{Tabla de direccionamiento IP.}
\end{table}

Los switches S1, S2, S3 y S4 son dispositivos de capa 2, por lo que no
necesitan una dirección IP para realizar su función de conmutación. Si se
requiere administración remota de los switches, se puede asignar una
dirección IP a una interfaz VLAN de administración dentro de la subred
correspondiente.

Las subredes 7 y 8 quedan reservadas para futuras ampliaciones de la red.

\end{document}
